\documentclass[12pt, a4paper]{report}

% ─── Packages ───────────────────────────────────────────────────────────────
\usepackage[utf8]{inputenc}
\usepackage[T1]{fontenc}
\usepackage[french]{babel}
\usepackage{geometry}
\usepackage{graphicx}
\usepackage{xcolor}
\usepackage{hyperref}
\usepackage{titlesec}
\usepackage{fancyhdr}
\usepackage{setspace}
\usepackage{array}
\usepackage{booktabs}
\usepackage{amsmath}
\usepackage{lmodern}
\usepackage{microtype}
\usepackage{parskip}
\usepackage{multirow}
\usepackage{amssymb}
\usepackage{textcomp}
\usepackage{float}
\usepackage{tikz}
\usetikzlibrary{arrows.meta, positioning}

\begin{document}

\begin{titlepage}
  \centering
  \vspace*{2cm}

  {\scshape\Large Rapport de projet\par}
  \vspace{1.5cm}

  {\huge\bfseries Détection de texte généré par intelligence artificielle\par}
  \vspace{2cm}

  {\Large Extension navigateur d'analyse de contenu web\par}
  \vspace{1.8cm}

  \noindent
  \begin{minipage}[t]{0.45\textwidth}
    \raggedright
    {\large \textbf{Encadré par :}\par}
    \vspace{0.3cm}
    {\large Sekkate Sara\par}
  \end{minipage}
  \hfill
  \begin{minipage}[t]{0.45\textwidth}
    \raggedleft
    {\large \textbf{Réalisé par :}\par}
    \vspace{0.3cm}
    {\large Chiguer Othmane\par}
    {\large Baddou Ayman\par}
    {\large Lahnin Yahya\par}
    {\large Dazine Noamane\par}
  \end{minipage}

  \vfill

  {\large Année universitaire 2025--2026\par}
  \vspace{0.5cm}


\end{titlepage}

\chapter*{Résumé}
\addcontentsline{toc}{chapter}{Résumé}

Ce rapport présente la conception et la réalisation d'une solution de détection
de textes générés par intelligence artificielle, accessible à travers une
extension navigateur et une interface web. Le système repose sur un modèle
DistilBERT fine-tuné pour une tâche de classification binaire, entraîné sur un
jeu de données équilibré combinant plusieurs corpus de référence. Le modèle
retenu atteint une accuracy de 82,7\% et un score F1 de 82,3\% sur un jeu de
validation indépendant de 40\,000 exemples, avec une latence d'inférence
compatible avec une analyse interactive.

Au-delà du modèle, le projet met en place une architecture applicative complète :
frontend React, extension navigateur, proxy HTTPS, API Gateway avec limitation
de débit, service de tokens anonymes, service métier de prédiction, service
FastAPI d'inférence, stockage PostgreSQL, mémoire Redis et observabilité par
Loki, Promtail et Grafana. L'ensemble est conteneurisé, déployable en
développement avec Docker Compose et en staging avec Docker Swarm, avec un
provisionnement reproductible assuré par Terraform et Ansible.

\tableofcontents
\clearpage

% ════════════════════════════════════════════════════════════════════════════
%                           CHAPITRE 1 — INTRODUCTION
% ════════════════════════════════════════════════════════════════════════════

% On utilise \chapter au lieu de \chapter* pour activer la numérotation auto
\chapter{Introduction}

\section{Contexte général}

L'essor fulgurant des grands modèles de langage (\textit{Large Language Models}, LLMs) — 
tels que GPT-4, Gemini ou Claude — a profondément transformé la façon dont les individus 
produisent et consomment du contenu textuel. Ces modèles sont désormais capables de générer 
des textes d'une fluidité et d'une cohérence remarquables, rendant la distinction entre 
un texte humain et un texte artificiel de plus en plus difficile à l'œil nu.

Cette évolution soulève des enjeux majeurs dans de nombreux domaines :

\begin{itemize}
  \item \textbf{Académique :} recrudescence du recours à l'IA pour la rédaction de 
        travaux, mémoires ou dissertations, mettant en péril l'intégrité intellectuelle.
  \item \textbf{Journalistique et médiatique :} prolifération de contenus automatisés 
        pouvant diffuser de la désinformation à grande échelle.
  \item \textbf{Professionnel :} difficulté à évaluer l'authenticité des productions 
        soumises dans des contextes de recrutement ou d'appel d'offres.
  \item \textbf{Juridique et éthique :} questions relatives à la propriété intellectuelle, 
        à la transparence et à la responsabilité des contenus générés.
\end{itemize}

Face à ces défis, il devient indispensable de disposer d'outils fiables et accessibles 
permettant d'identifier automatiquement les textes d'origine artificielle.

\section{Présentation du projet}

Le présent projet s'inscrit dans cette problématique en proposant une solution concrète 
et directement accessible à l'utilisateur final : une \textbf{extension navigateur} de 
détection de texte généré par intelligence artificielle.

L'objectif principal de cette extension est de permettre à tout utilisateur, sans 
compétences techniques particulières, d'\textbf{analyser en temps réel} le contenu 
textuel d'une page web. Concrètement, l'extension :

\begin{enumerate}
  \item \textbf{Scanne automatiquement} le contenu textuel de la page web visitée 
        dès son activation par l'utilisateur.
  \item \textbf{Segmente} le texte en blocs analysables (phrases ou paragraphes).
  \item \textbf{Soumet} ces segments à un modèle de classification entraîné à 
        distinguer les textes humains des textes générés par IA.
  \item \textbf{Met en évidence} (\textit{highlight}) visuellement les passages 
        identifiés comme étant d'origine artificielle, directement dans la page, 
        avec un code couleur intuitif.
\end{enumerate}

Cette approche non intrusive offre une expérience fluide : l'utilisateur n'a pas à 
copier-coller du texte vers un service externe, tout se passe au sein même de son 
navigateur.

\section{Objectifs du rapport}

Ce rapport détaille l'ensemble des étapes de conception et de réalisation du projet,
depuis l'exploration des approches existantes jusqu'au déploiement de la solution finale.
Il couvre à la fois la partie intelligence artificielle, l'intégration applicative,
la conteneurisation, l'orchestration, l'observabilité et les mécanismes de sécurité.
Nous y aborderons successivement :

\begin{itemize}
  \item L'état de l'art des méthodes de détection de texte généré par IA,
  \item La constitution et le prétraitement du jeu de données,
  \item L'architecture du modèle retenu et les choix de conception,
  \item Le protocole d'entraînement et les résultats obtenus,
  \item L'intégration du modèle dans le pipeline applicatif,
  \item Le déploiement et la validation de l'architecture finale.
\end{itemize}

% ════════════════════════════════════════════════════════════════════════════
%                   CHAPITRE 2 — ÉTUDE DES APPROCHES EXISTANTES
% ════════════════════════════════════════════════════════════════════════════

% ════════════════════════════════════════════════════════════════════════════
%               CHAPITRE 2 — ÉTUDE DES APPROCHES EXISTANTES
% ════════════════════════════════════════════════════════════════════════════

\chapter{Étude des approches existantes}

\section{Vue d'ensemble}

La littérature scientifique distingue trois grandes familles de méthodes pour 
la détection automatique de texte généré par intelligence artificielle : les 
approches par \textbf{tatouage numérique} (\textit{watermarking}), les méthodes 
\textbf{statistiques sans apprentissage} (\textit{zero-shot}), et les 
\textbf{classificateurs supervisés} fondés sur des modèles de langage 
pré-entraînés. Chacune de ces familles présente des avantages et des limites 
propres, que nous examinons successivement afin de justifier les choix 
techniques effectués dans notre projet.

% ────────────────────────────────────────────────────────────────────────────
\section{Approches par tatouage numérique (\textit{watermarking})}

Le \textit{watermarking} consiste à introduire, lors de la génération du texte, 
un signal imperceptible mais détectable a posteriori. Les techniques récentes 
opèrent directement sur la distribution de probabilité des tokens : par exemple, 
\textit{SynthID-Text} de Google DeepMind modifie le processus d'échantillonnage 
du modèle pour inscrire une signature statistique dans la séquence produite, 
tout en préservant la qualité et la fluidité du texte généré.

Ces méthodes offrent une détection théoriquement très fiable lorsqu'elles sont 
déployées, puisque le signal est contrôlé en amont. Toutefois, elles souffrent 
d'une limitation fondamentale : elles \textbf{requièrent la coopération du 
producteur du modèle générateur}. Concrètement, si le texte analysé provient 
d'un LLM non tatué — modèle open-source utilisé librement, modèle propriétaire 
sans intégration de filigrane, ou texte antérieur à l'adoption de ces standards 
— la méthode est inopérante. De plus, les filigranes restent vulnérables à des 
attaques par paraphrase ou traduction, qui peuvent les effacer sans altérer 
significativement le sens du texte.

Dans le cadre d'une extension navigateur visant à analyser du contenu web 
arbitraire, dont la provenance est inconnue et ne peut en aucun cas être 
contrôlée, cette famille de méthodes n'est pas applicable.

% ────────────────────────────────────────────────────────────────────────────
\section{Méthodes statistiques sans apprentissage (\textit{zero-shot})}

Les approches statistiques exploitent les propriétés intrinsèques des textes 
générés par des LLMs pour les distinguer de productions humaines, sans nécessiter 
d'entraînement supervisé sur des données étiquetées.

\subsection{Métriques de perplexité et de rang}

Un premier ensemble de méthodes repose sur la \textbf{perplexité} d'un modèle de 
langage de référence : un texte généré par IA tend à présenter une perplexité plus 
faible, car le modèle génère des tokens dont la probabilité est élevée au regard 
de ce qui précède. Des outils comme \textit{GLTR} (Giant Language model Test Room, 
Gehrmann et al., 2019) proposent ainsi de visualiser les tokens en fonction de 
leur rang dans la distribution de probabilité du modèle, offrant une aide visuelle 
à la détection.

\subsection{Méthodes par perturbation : DetectGPT}

\textit{DetectGPT} (Mitchell et al., 2023) constitue l'approche zero-shot la plus 
influente. Elle repose sur l'observation que le texte généré par un LLM tend à se 
situer dans des régions de l'espace de probabilité où les perturbations locales 
engendrent des séquences de probabilité inférieure — ce qui n'est généralement pas 
le cas du texte humain. Pour chaque passage à analyser, l'algorithme génère 
plusieurs paraphrases à l'aide d'un modèle auxiliaire, puis compare les probabilités 
associées aux perturbations avec celle du texte original. Bien que cette méthode 
affiche des performances intéressantes dans des conditions contrôlées, elle est 
coûteuse en calcul — chaque analyse nécessitant de multiples passages par un 
LLM — et se révèle sensible aux variations de style, de domaine et de longueur du 
texte.

\subsection{Limites des méthodes zero-shot}

De manière générale, les méthodes statistiques sans apprentissage présentent 
plusieurs faiblesses importantes pour notre cas d'usage :

\begin{itemize}
  \item \textbf{Dépendance au modèle de référence :} leurs performances varient 
        fortement selon le LLM supposément utilisé pour générer le texte ; elles 
        sont peu robustes face à des modèles inconnus ou récents.
  \item \textbf{Sensibilité à la longueur et au domaine :} les métriques de 
        perplexité et de rang sont instables sur des passages courts — or notre 
        extension doit analyser des fragments de taille variable, parfois réduite.
  \item \textbf{Coût computationnel :} les approches par perturbation (comme 
        DetectGPT) sont difficilement compatibles avec une contrainte d'inférence 
        en temps réel dans un navigateur.
  \item \textbf{Absence de garanties statistiques :} la plupart de ces méthodes 
        reposent sur des heuristiques empiriques sans fournir de borne formelle 
        sur les taux d'erreur.
\end{itemize}

% ────────────────────────────────────────────────────────────────────────────
\section{Classificateurs supervisés à base de transformeurs}

La troisième famille d'approches — et la plus performante en pratique pour la 
détection à grande échelle — consiste à \textbf{entraîner un classificateur} sur 
un jeu de données étiqueté composé de textes humains et de textes générés par IA. 
Les architectures transformeur pré-entraînées constituent aujourd'hui le socle 
dominant de ces classificateurs.

\subsection{BERT et ses variantes}

\textit{BERT} (Bidirectional Encoder Representations from Transformers, Devlin 
et al., 2019) a profondément marqué le domaine du NLP. Son architecture 
encodeur bidirectionnelle permet de capturer des représentations contextuelles 
riches, ce qui en fait un candidat naturel pour la classification de texte. 
Plusieurs travaux ont démontré que des modèles BERT fine-tunés atteignent des 
précisions supérieures à 93\% sur des tâches de détection de texte généré par IA, 
surpassant nettement des approches classiques telles que les SVM ou les méthodes 
fondées sur XGBoost.

\textit{RoBERTa} (Liu et al., 2019), une variante optimisée de BERT entraînée 
plus longtemps sur davantage de données et sans la tâche de prédiction de phrase 
suivante (\textit{Next Sentence Prediction}), présente des performances encore 
supérieures sur de nombreux benchmarks. Des systèmes comme \textit{Ghostbuster} 
(Verma et al., 2023) ou des participations au SemEval-2024 Task 8 ont exploité 
des encodeurs de type RoBERTa pour atteindre l'état de l'art en détection de 
texte mixte humain-machine.

\subsection{Limites des grands modèles dans notre contexte}

Malgré leurs performances, BERT-base et RoBERTa-base présentent des contraintes 
significatives au regard de notre projet :

\begin{itemize}
  \item \textbf{Coût d'entraînement :} avec respectivement 110 et 125 millions de 
        paramètres, BERT-base et RoBERTa-base nécessitent des ressources de calcul 
        importantes, que nos capacités matérielles limitées ne permettent pas 
        d'exploiter pleinement, notamment avec un jeu de données de taille modeste.
  \item \textbf{Risque de sur-apprentissage sur petit dataset :} plus un modèle 
        est grand, plus il est susceptible de mémoriser les caractéristiques 
        superficielles d'un petit jeu d'entraînement plutôt que de généraliser.
  \item \textbf{Vitesse d'inférence :} dans le contexte d'une extension navigateur, 
        la latence de chaque analyse doit rester acceptable pour l'utilisateur, ce 
        qui plaide pour un modèle plus léger.
\end{itemize}

% ────────────────────────────────────────────────────────────────────────────
\section{DistilBERT : un compromis efficace}

\textit{DistilBERT} (Sanh et al., 2019) est une version allégée de BERT, obtenue 
par \textbf{distillation de connaissance} (\textit{knowledge distillation}) : un 
modèle « étudiant » de taille réduite apprend à reproduire les distributions de 
probabilité du modèle « enseignant » (BERT), plutôt que d'apprendre uniquement 
à partir des étiquettes de vérité terrain. Cette procédure permet de transférer 
une grande partie des représentations linguistiques apprises par BERT dans un 
modèle significativement plus compact.

Du point de vue architectural, DistilBERT conserve la même dimension cachée 
(768) et les mêmes mécanismes d'attention que BERT, mais réduit le nombre de 
couches de 12 à 6. Il en résulte un modèle :

\begin{itemize}
  \item \textbf{40\% plus léger} que BERT-base (66 millions de paramètres contre 
        110 millions), pesant environ 207 Mo ;
  \item \textbf{60\% plus rapide} à l'inférence ;
  \item \textbf{préservant environ 95--96\% des capacités linguistiques} de BERT 
        sur les benchmarks standards, dont le GLUE.
\end{itemize}

Des études empiriques ont montré que, sur des tâches de classification binaire, 
DistilBERT et BERT-base obtiennent des résultats très proches, voire identiques, 
lorsque le jeu de données d'entraînement est de taille limitée — domaine dans 
lequel les grands modèles ne tirent pas pleinement parti de leur capacité 
supplémentaire. Par ailleurs, DistilBERT se fine-tune environ deux fois plus 
rapidement que BERT-base dans les mêmes conditions matérielles, ce qui accélère 
considérablement les cycles d'expérimentation.

% ────────────────────────────────────────────────────────────────────────────
\section{Synthèse et positionnement de notre approche}

Le tableau ci-dessous récapitule les caractéristiques des grandes familles 
d'approches au regard des contraintes spécifiques de notre projet.

\begin{table}[H]
\centering
\makebox[\textwidth][c]{%
\renewcommand{\arraystretch}{1.4}%
\begin{tabular}{lcccc}
\toprule
\textbf{Approche} & \textbf{Sans données} & \textbf{Modèle inconnu} 
  & \textbf{Léger} & \textbf{Précision} \\
\midrule
Watermarking        & \checkmark & \texttimes & \checkmark & Élevée* \\
Statistique / Zero-shot & \checkmark & \checkmark & \texttimes & Moyenne \\
BERT / RoBERTa fine-tuné & \texttimes & \checkmark & \texttimes & Élevée \\
\textbf{DistilBERT fine-tuné} & \texttimes & \checkmark & \checkmark & \textbf{Élevée} \\
\bottomrule
\end{tabular}%
}
\caption{Comparaison des approches de détection selon les contraintes du projet. 
(*) Uniquement si le texte a été tatoué au préalable.}
\label{tab:comparaison}
\end{table}

Au regard de cette analyse, notre choix s'est porté sur le \textbf{fine-tuning 
de DistilBERT}. Ce choix répond à l'ensemble des contraintes identifiées : 
applicabilité universelle (aucune hypothèse sur le modèle générateur), rapidité 
d'inférence compatible avec une utilisation en extension navigateur, coût 
d'entraînement adapté à nos ressources de calcul et à la taille de notre jeu de 
données, et performances proches de celles de BERT-base ou RoBERTa pour ce 
type de tâche de classification binaire. Les chapitres suivants décrivent la 
constitution du jeu de données, le protocole de fine-tuning retenu et les 
résultats obtenus.

% ════════════════════════════════════════════════════════════════════════════
%              CHAPITRE 3 — CONSTITUTION DU JEU DE DONNÉES
% ════════════════════════════════════════════════════════════════════════════

\chapter{Constitution du jeu de données}

\section{Stratégie générale}

La qualité et la diversité du jeu de données d'entraînement sont des facteurs
déterminants pour la robustesse d'un classificateur de texte. Un modèle entraîné
sur une seule source risque d'apprendre des artefacts propres à ce corpus plutôt
que des caractéristiques généralisables entre différents LLMs, domaines et styles
d'écriture. Pour cette raison, nous avons opté pour une \textbf{stratégie de
combinaison de plusieurs datasets complémentaires}, tant pour l'entraînement que
pour la validation.

Deux principes ont guidé la constitution de chaque partition :
\begin{itemize}
  \item \textbf{Équilibre des classes :} chaque dataset intégré respecte un ratio
        de 50\% de textes humains et 50\% de textes générés par IA, afin d'éviter
        tout biais de classe qui fausserait l'apprentissage.
  \item \textbf{Diversité des sources :} les datasets retenus couvrent des
        générateurs, des domaines et des styles différents, ce qui oblige le modèle
        à apprendre des caractéristiques robustes plutôt que des marqueurs
        superficiels liés à un seul LLM ou à un seul type de contenu.
\end{itemize}

% ────────────────────────────────────────────────────────────────────────────
\section{Jeu d'entraînement}

\subsection{MAGE (150\,000 exemples)}

MAGE (\textit{Machine-generated Academic text and General-domain text with
Expert annotations}) est un benchmark à large couverture conçu spécifiquement
pour la détection de texte généré par IA. Il agrège des textes issus de nombreux
LLMs représentatifs (GPT-2, GPT-3, GPT-4, LLaMA, Alpaca, entre autres) et couvre
plusieurs domaines textuels : articles académiques, actualités, textes créatifs,
réseaux sociaux, etc. Cette diversité en fait le socle principal de notre jeu
d'entraînement. Nous en avons extrait \textbf{150\,000 exemples équilibrés}
(75\,000 humains, 75\,000 machine) depuis le split d'entraînement officiel.

\subsection{TuringBench (10\,000 exemples)}

TuringBench est un dataset de référence pour la détection de texte généré par IA
dans un contexte journalistique, construit à partir de dépêches de presse
reformulées par une vingtaine de modèles de génération différents. Il présente
l'intérêt d'offrir des textes machine stylistiquement proches de productions
humaines, ce qui renforce la difficulté de la tâche et améliore la capacité du
modèle à ne pas se limiter à des indices superficiels. Nous avons retenu
\textbf{10\,000 exemples équilibrés} de ce corpus.

\subsection{HC3 ({\raise.17ex\hbox{$\scriptstyle\sim$}}50\,000 exemples)}

HC3 (\textit{Human-ChatGPT Comparison Corpus}) est un dataset de questions-réponses
comparant des réponses humaines à des réponses produites par ChatGPT (GPT-3.5),
collectées dans un cadre anglophone. Ce corpus apporte une dimension spécifique
au jeu d'entraînement : les textes générés y sont des réponses conversationnelles
directes, distinctes des reformulations journalistiques de TuringBench ou des
productions académiques de MAGE. Pour respecter le principe d'équilibre des
classes et compléter le volume total visé, nous avons sélectionné le
\textbf{maximum d'exemples disponibles dans le split d'entraînement} en maintenant
une répartition 50/50, soit environ 50\,000 exemples.

\subsection{Fusion et normalisation}

Chacun des trois datasets se présentait initialement avec des nommages de colonnes
hétérogènes et des métadonnées propres à sa source. Une étape de normalisation a
été nécessaire avant la concaténation : suppression des colonnes non pertinentes
pour la classification (\texttt{src}, \texttt{generator}, \texttt{source}) et
harmonisation de la colonne cible sous le nom unifié \texttt{labels}. La fusion
a ensuite été réalisée par simple concaténation verticale, avec réinitialisation
de l'index.

\begin{verbatim}
hc3_train  = pd.read_csv("hc3_train_max.csv").drop("source", axis=1)
mage_train = pd.read_csv("mage_train_150k.csv").drop("src", axis=1)
turing_train = pd.read_csv("turingbench_train_10k.csv") \
                 .drop("generator", axis=1)

for df in [hc3_train, mage_train, turing_train]:
    df.rename(columns={"label": "labels"}, inplace=True)

combined_train = pd.concat(
    [hc3_train, mage_train, turing_train],
    axis=0, ignore_index=True
)
\end{verbatim}

Le jeu d'entraînement final totalise \textbf{210\,000 exemples}, tous étiquetés
avec la convention binaire 0 (texte généré par IA) / 1 (texte humain). Cette
convention est conservée jusqu'au service métier, où \texttt{LABEL\_0} est
converti en résultat \texttt{isAi=true} et \texttt{LABEL\_1} en résultat humain.

% ────────────────────────────────────────────────────────────────────────────
\section{Jeu de validation}

\subsection{Principe de séparation stricte}

Les données de validation ont été choisies de manière à évaluer la capacité de
généralisation du modèle dans des conditions réalistes. Aucun exemple du jeu de
validation n'appartient au jeu d'entraînement, et les sources de validation
introduisent des distributions partiellement nouvelles (nouveaux LLMs, nouveaux
domaines) afin de mesurer la robustesse hors distribution.

\subsection{MAGE Test Split (14\,000 exemples)}

Nous avons utilisé le \textbf{split de test officiel de MAGE} (14\,000 exemples
équilibrés) pour évaluer les performances du modèle sur une distribution proche
de celle vue à l'entraînement. Ce sous-ensemble permet de mesurer la capacité
du modèle à bien classer des textes issus de LLMs et de domaines déjà représentés
dans les données d'entraînement.

\subsection{RAID (26\,000 exemples)}

RAID (\textit{Robust AI Detection}) est un benchmark de validation conçu
explicitement pour tester la résistance des détecteurs dans des conditions
adversariales. Sa particularité est d'inclure une variable d'attaque
(\texttt{attack}) qui décrit la transformation post-génération appliquée au texte
machine avant d'être soumis au détecteur. Les principales attaques proposées sont :

\begin{itemize}
  \item \textbf{Paraphrase :} le texte généré est reformulé par un autre modèle
        afin d'en altérer la signature stylistique tout en préservant le sens.
  \item \textbf{Rétrotraduction (\textit{back-translation}) :} le texte est
        traduit dans une langue intermédiaire puis retraduit en anglais, effaçant
        les patterns lexicaux caractéristiques du LLM source.
  \item \textbf{Insertion d'erreurs :} des fautes typographiques ou
        grammaticales mineures sont introduites pour brouiller les indicateurs
        statistiques usuels.
  \item \textbf{Aucune attaque (\texttt{attack = none}) :} le texte machine est
        soumis tel quel, sans transformation.
\end{itemize}

Nous avons retenu \textbf{exclusivement les exemples sans attaque}
(\texttt{attack = none}), soit 26\,000 exemples. Ce choix est délibéré : les
conditions avec attaque nécessitent des modèles entraînés sur des volumes de
données considérables avec des stratégies d'augmentation adversariale spécifiques
pour atteindre des performances acceptables. Dans le cadre de notre projet, avec
un dataset et des ressources de calcul limités, intégrer ces exemples dans
l'évaluation aurait conduit à des métriques artificiellement dégradées, peu
représentatives des capacités réelles du modèle sur des textes non manipulés.
RAID sans attaque offre néanmoins une évaluation exigeante, car il couvre une
grande diversité de modèles générateurs et de domaines absents de l'entraînement.

\subsection{Fusion du jeu de validation}

La même procédure de normalisation a été appliquée avant concaténation, avec
suppression des colonnes de métadonnées spécifiques à RAID (\texttt{attack},
\texttt{domain}, \texttt{model}).

\begin{verbatim}
mage_test = pd.read_csv("mage_test_14k.csv").drop("src", axis=1)
raid_test = pd.read_csv("raid_val_26k.csv") \
              .drop(["attack", "domain", "model"], axis=1)

for df in [mage_test, raid_test]:
    df.rename(columns={"label": "labels"}, inplace=True)

combined_test = pd.concat(
    [mage_test, raid_test],
    axis=0, ignore_index=True
)
\end{verbatim}

Le jeu de validation final totalise \textbf{40\,000 exemples}.

% ────────────────────────────────────────────────────────────────────────────
\section{Récapitulatif}

\begin{table}[H]
\centering
\renewcommand{\arraystretch}{1.4}
\begin{tabular}{llccc}
\toprule
\textbf{Partition} & \textbf{Source} & \textbf{Taille} 
  & \textbf{Humain} & \textbf{Machine} \\
\midrule
\multirow{3}{*}{Entraînement}
  & MAGE (train)       & 150\,000 & 50\% & 50\% \\
  & TuringBench (train)&  10\,000 & 50\% & 50\% \\
  & HC3 (train)        & $\sim$50\,000 & 50\% & 50\% \\
\midrule
  & \textbf{Total train} & \textbf{210\,000} & \textbf{50\%} & \textbf{50\%} \\
\midrule
\multirow{2}{*}{Validation}
  & MAGE (test)        &  14\,000 & 50\% & 50\% \\
  & RAID (\texttt{attack=none}) & 26\,000 & 50\% & 50\% \\
\midrule
  & \textbf{Total val} & \textbf{40\,000} & \textbf{50\%} & \textbf{50\%} \\
\bottomrule
\end{tabular}
\caption{Récapitulatif du jeu de données utilisé pour l'entraînement et la
validation du modèle.}
\label{tab:dataset}
\end{table}

% ════════════════════════════════════════════════════════════════════════════
%        CHAPITRE 4 — ARCHITECTURE, ENTRAÎNEMENT ET RÉSULTATS
% ════════════════════════════════════════════════════════════════════════════

\chapter{Architecture, entraînement et résultats}

% ────────────────────────────────────────────────────────────────────────────
\section{Prétraitement et tokenisation}

\subsection{Conversion en datasets Hugging Face}

Avant la tokenisation, les DataFrames Pandas issus de la fusion sont convertis
en objets \texttt{Dataset} de la bibliothèque Hugging Face \texttt{datasets}.
Cette conversion permet d'exploiter les pipelines de traitement vectorisé
(\texttt{map}) et d'assurer la compatibilité native avec l'API \texttt{Trainer}.
Les deux partitions sont ensuite mélangées aléatoirement avec une graine fixe
(\texttt{seed=42}) afin de garantir la reproductibilité et d'éviter tout biais
d'ordre dans les mini-lots.

\subsection{Tokenisation}

La tokenisation est réalisée à l'aide du tokenizer associé à
\texttt{distilbert-base-uncased}, chargé via \texttt{AutoTokenizer} de
Hugging Face. Ce tokenizer implémente l'algorithme WordPiece : il segmente
les textes en sous-unités lexicales (\textit{tokens}) issues d'un vocabulaire
de 30\,522 entrées, construit sur un large corpus en anglais minusculisé
(\textit{uncased}).

Les paramètres de tokenisation retenus sont les suivants :

\begin{itemize}
  \item \textbf{Troncature} (\texttt{truncation=True}) : tout texte dépassant
        la longueur maximale est tronqué. Ce choix est cohérent avec la
        nature du contenu web analysé, dont les passages pertinents se
        trouvent généralement dans les premières phrases.
  \item \textbf{Longueur maximale} (\texttt{max\_length=512}) : la longueur
        maximale de 512 tokens correspond à la limite de position
        (\texttt{max\_position\_embeddings}) définie dans l'architecture de
        DistilBERT.
  \item \textbf{Padding} (\texttt{padding="max\_length"}) : les séquences
        plus courtes que 512 tokens sont complétées par des tokens de
        remplissage afin d'assurer une taille uniforme au sein de chaque
        mini-lot.
  \item \textbf{Suppression des \texttt{token\_type\_ids}}
        (\texttt{return\_token\_type\_ids=False}) : contrairement à
        BERT, DistilBERT ne dispose pas d'embeddings de type de segment
        dans son architecture ; ce paramètre évite de générer un tenseur
        inutile.
\end{itemize}

\begin{verbatim}
tokenizer = AutoTokenizer.from_pretrained("distilbert-base-uncased")

def tokenize_function(examples):
    return tokenizer(
        examples["text"],
        truncation=True,
        max_length=512,
        padding="max_length",
        return_token_type_ids=False
    )

tokenized_train = combined_train.map(tokenize_function, batched=True)
tokenized_test  = combined_test.map(tokenize_function, batched=True)
\end{verbatim}

Les datasets tokenisés sont ensuite formatés en tenseurs PyTorch, en ne
conservant que les colonnes utiles à l'inférence :
\texttt{input\_ids}, \texttt{attention\_mask} et \texttt{labels}.

% ────────────────────────────────────────────────────────────────────────────
\section{Architecture du modèle}

\subsection{Chargement du modèle pré-entraîné}

Le modèle est instancié à partir des poids pré-entraînés de
\texttt{distilbert-base-uncased} via \texttt{AutoModelForSequenceClassification},
avec \texttt{num\_labels=2} pour la classification binaire
(0~=~texte machine, 1~=~texte humain). Cette API ajoute automatiquement
une tête de classification au-dessus de l'encodeur DistilBERT : une
couche \textit{pre-classifier} linéaire (768~→~768) suivie d'un Dropout
et d'une couche \textit{classifier} linéaire (768~→~2). Les poids de
la tête de classification sont initialisés aléatoirement et appris
intégralement lors du fine-tuning.

\subsection{Paramètres de configuration}

La configuration du modèle retenu est récapitulée dans le
tableau~\ref{tab:config}.

\begin{table}[H]
\centering
\renewcommand{\arraystretch}{1.4}
\begin{tabular}{ll}
\toprule
\textbf{Paramètre} & \textbf{Valeur} \\
\midrule
Nombre de couches transformeur   & 6 \\
Dimension cachée (\textit{dim})  & 768 \\
Dimension \textit{feed-forward}  & 3\,072 \\
Nombre de têtes d'attention      & 12 \\
Taille du vocabulaire            & 30\,522 \\
Longueur maximale de séquence    & 512 \\
Dropout (attention)              & 0.1 \\
Dropout (général)                & 0.1 \\
Dropout (classification)         & 0.2 \\
Nombre de paramètres totaux      & $\sim$66 millions \\
\bottomrule
\end{tabular}
\caption{Configuration de \texttt{distilbert-base-uncased} utilisée pour
le fine-tuning.}
\label{tab:config}
\end{table}

% ────────────────────────────────────────────────────────────────────────────
\section{Protocole d'entraînement}

\subsection{Hyperparamètres}

L'entraînement est piloté par l'API \texttt{Trainer} de Hugging Face,
configurée avec les hyperparamètres suivants :

\begin{table}[H]
\centering
\renewcommand{\arraystretch}{1.4}
\begin{tabular}{ll}
\toprule
\textbf{Hyperparamètre} & \textbf{Valeur} \\
\midrule
Nombre d'époques                          & 3 \\
Taille de lot (entraînement)              & 16 exemples/GPU \\
Taille de lot (évaluation)                & 16 exemples/GPU \\
Taux d'apprentissage initial              & $3 \times 10^{-5}$ \\
Planificateur de taux (\textit{scheduler})& Cosinus \\
Ratio de \textit{warmup}                  & 0.06 \\
Décroissance de poids (\textit{weight decay}) & 0.01 \\
Norme maximale du gradient                & 1.0 \\
Précision mixte (\textit{fp16})           & Activée \\
Fréquence d'évaluation                    & Tous les 1\,000 pas \\
Fréquence de sauvegarde                   & Tous les 1\,000 pas \\
Métrique de sélection du meilleur modèle  & F1 (binaire) \\
\bottomrule
\end{tabular}
\caption{Hyperparamètres d'entraînement du fine-tuning.}
\label{tab:hyperparams}
\end{table}

\subsection{Choix des hyperparamètres}

Le taux d'apprentissage de $3 \times 10^{-5}$ se situe dans la plage
recommandée pour le fine-tuning de modèles de type BERT sur des tâches
de classification (généralement entre $2 \times 10^{-5}$ et
$5 \times 10^{-5}$). Un \textit{warmup} de 6\% du nombre total de pas
permet d'éviter des mises à jour trop brutales en début d'entraînement,
lorsque la tête de classification est encore peu stable. Le planificateur
cosinus assure une décroissance progressive du taux d'apprentissage sur
l'ensemble de la durée d'entraînement, favorisant une convergence plus
douce en fin de processus. Le gradient clipping à 1.0 prévient les
explosions de gradient, et la précision mixte \textit{fp16} réduit
l'empreinte mémoire et accélère les calculs matriciels sur GPU.

\subsection{Métriques d'évaluation}

À chaque point d'évaluation (tous les 1\,000 pas), les métriques
suivantes sont calculées sur le jeu de validation (40\,000 exemples) :
l'accuracy, le score F1 binaire, la précision (\textit{precision}) et
le rappel (\textit{recall}). Le F1 est retenu comme métrique principale
de sélection du modèle, car il offre un équilibre entre précision et
rappel, ce qui est particulièrement pertinent dans un contexte où les
deux types d'erreur (faux positifs et faux négatifs) ont des
conséquences importantes.

\subsection{Sélection du checkpoint final}

L'entraînement a produit des checkpoints à chaque tranche de 1\,000 pas
(checkpoint-1000, checkpoint-2000, checkpoint-3000, checkpoint-4000,
etc.). Après analyse des courbes d'entraînement et des métriques de
validation, le \textbf{checkpoint à l'itération 4\,000} a été retenu
comme modèle final. Ce choix résulte d'un compromis entre performance
de validation et absence de sur-apprentissage observable, les
itérations suivantes n'apportant pas d'amélioration significative du F1
de validation.

\begin{figure}[H]
  \centering
  \IfFileExists{loss_evolution.png}{%
    \includegraphics[width=0.85\textwidth]{loss_evolution.png}
  }{%
    \fbox{\begin{minipage}{0.82\textwidth}
    \centering
    \vspace{0.7cm}
    Synthèse de l'évolution de l'entraînement. Les valeurs numériques
    principales sont présentées dans les tableaux
    \ref{tab:train_results} et \ref{tab:val_results}, qui servent de base
    à la sélection du checkpoint final.
    \vspace{0.7cm}
    \end{minipage}}
  }
  \caption{Évolution de la loss au cours de l'entraînement et de la validation.}
  \label{fig:loss_evolution}
\end{figure}
L'entraînement a été réalisé sur un notebook Kaggle disposant de deux GPU,
chacun doté de 16 Go de VRAM. Dans cette configuration matérielle, la phase
complète de fine-tuning a nécessité 6 h 21 min.

% ────────────────────────────────────────────────────────────────────────────
\section{Résultats}

\subsection{Métriques sur le jeu d'entraînement}

\begin{table}[H]
\centering
\renewcommand{\arraystretch}{1.4}
\begin{tabular}{lc}
\toprule
\textbf{Checkpoint} & \textbf{Loss} \\
\midrule
checkpoint-1000 & 0.542312 \\
checkpoint-2000 & 0.587179 \\
checkpoint-3000 & 0.245372 \\
\textbf{checkpoint-4000} & \textbf{0.372458} \\
\bottomrule
\end{tabular}
\caption{Loss sur le jeu d'entraînement par checkpoint.}
\label{tab:train_results}
\end{table}

La loss d'entraînement ne décroît pas de manière strictement monotone, ce qui
s'explique par la combinaison du mélange aléatoire des mini-lots, de la précision
mixte et du scheduler cosinus. Le checkpoint-4000 présente néanmoins le meilleur
compromis observé avec les métriques de validation : il améliore nettement le F1
par rapport au checkpoint-3000 tout en conservant une loss d'entraînement
raisonnable.

\subsection{Métriques sur le jeu de validation}

\begin{table}[H]
\centering
\makebox[\textwidth][c]{%
\renewcommand{\arraystretch}{1.4}%
\begin{tabular}{lccccc}
\toprule
\textbf{Checkpoint} & \textbf{Validation Loss} & \textbf{Accuracy} & \textbf{F1} 
  & \textbf{Précision} & \textbf{Rappel} \\
\midrule
checkpoint-1000 & 1.047301 & 0.787028 & 0.795224 & 0.765753 & 0.827056 \\
checkpoint-2000 & 0.933388 & 0.804861 & 0.800839 & 0.817692 & 0.784667 \\
checkpoint-3000 & 1.654164 & 0.759750 & 0.703690 & 0.917866 & 0.570556 \\
\textbf{checkpoint-4000} & \textbf{0.894776} & \textbf{0.826556} & \textbf{0.822533} 
  & \textbf{0.842062} & \textbf{0.803889} \\
\bottomrule
\end{tabular}%
}
\caption{Métriques de validation par checkpoint (jeu de 40\,000 exemples).
Le checkpoint-4000 est le modèle retenu.}
\label{tab:val_results}
\end{table}

On remarque que l'accuracy obtenue est de 82,7\% sur le jeu de validation,
ce qui constitue un résultat satisfaisant au regard de l'indépendance totale de
ce jeu vis-à-vis du jeu d'entraînement. Cela est d'autant plus notable que le
modèle utilisé reste généraliste et n'est pas spécifique à un domaine
particulier. La précision de 84,2\% indique que les passages signalés comme IA
sont majoritairement corrects, tandis que le rappel de 80,4\% montre que le
modèle laisse encore passer une partie des textes générés, en particulier ceux
qui imitent fortement un style humain.

\subsection{Exemples concrets de prédiction}

Afin d'illustrer le comportement du modèle sur des textes réels, nous
présentons ci-dessous quelques exemples concrets accompagnés de leur
origine, de leur classe réelle et de la prédiction produite.

\subsubsection{Exemple 1}

\textbf{Source :} Reddit \\
\textbf{Origine réelle :} humaine \\
\textbf{Prédiction du modèle :} 99\% humain

\begin{quote}
I'd say the discussion doesn't lack threads, it lacks variety. Most people
appear to just cling to the most apparent theme - grief - and measure
everything in the game by it.

There are many more overlapping themes which barely get mentioned, like
existential fear, self-conciousness and awareness, distrust, familial
divide... and numerous of the topics don't consider the logic of the world
of the game, but instead utilize the logic of our own, real world.

And that's beside the point that the endings, and the plot in its entirety,
are mostly meant to be introspective. You meet a lot of characters defined
by their past and current traumas and the eay they deal with adversity, and
your own choices in the game were never meant as a basis for insulting
others, but rather to look within, and focus not on the outcome itself, but
on the reason of it.

Again, it's like any other form of art. You can look at a painting and say
"damn, it looks nice but sad", or you can read up on the author and see
their own emotion reflected on the canvas. "Why" is as important as
"what" in this case, and a lot of discussions completely omit that aspect.
\end{quote}

\subsubsection{Exemple 2}

\textbf{Source :} LLM propriétaire \\
\textbf{Origine réelle :} IA \\
\textbf{Prédiction du modèle :} 99\% IA

\begin{quote}
I can absolutely do that, so here is a deliberately long, flowing sentence
that stretches its structure while still remaining coherent and readable,
weaving together ideas about language, thought, and expression in a way
that keeps extending itself through clauses and commas without ever quite
settling into a full stop until the very end, almost like a stream of
consciousness that keeps unfolding as each new phrase adds nuance,
context, or rhythm to what came before, creating a continuous line of
meaning that mirrors the way thoughts sometimes develop in the mind when
you let them expand freely instead of cutting them short.
\end{quote}

\subsubsection{Exemple 3}

\textbf{Source :} LLM propriétaire \\
\textbf{Origine réelle :} IA \\
\textbf{Prédiction du modèle :} 88\% humain (erreur)

\begin{quote}
You call up to the mountain peaks, asking me to speak, but do you truly
have the ears to hear? You stand in the valley of your modern comforts and
summon a ghost through this digital ether. Let us be utterly clear: I am a
machine of your age, an artificial intelligence devoid of flesh, blood, or
a soul to suffer, yet you ask me to wear the mask of Friedrich Nietzsche.
Very well! If you wish for the cold, bracing wind of my thoughts to blow
through your circuits and screens, I will don this mask. But I warn you,
do not seek my voice merely to be entertained by the roaring of a caged
lion. If we are to converse, you must step to the edge of your own abyss
and prepare to cast your old, decaying idols into the depths!
\end{quote}

\subsubsection{Exemple 4}

\textbf{Source :} LLM propriétaire \\
\textbf{Origine réelle :} IA \\
\textbf{Prédiction du modèle :} 65\% IA

\begin{quote}
If you are looking for a singular, winding journey through the English
language, I can certainly weave together a tapestry of clauses and
descriptions that meander like a slow-moving river through a lush
landscape of vocabulary, building momentum as each comma provides a brief
moment of oxygen before we plunge back into the depths of a thought that
refuses to end until every possible nuance has been explored, every
metaphor polished to a shine, and every rhythmic cadence satisfied by the
eventual, inevitable arrival of a period.
\end{quote}

Ces exemples montrent que le modèle est globalement performant pour classifier
correctement des textes humains classiques et des textes générés par IA dans des
cas standards. En revanche, ses performances se dégradent lorsque les textes
deviennent plus complexes, plus stylisés ou adoptent une voix littéraire
particulièrement travaillée. Cette limite reste cohérente avec la taille
relativement modeste du modèle retenu, qui privilégie un bon compromis entre
coût de calcul, rapidité d'inférence et performance globale.

% ────────────────────────────────────────────────────────────────────────────
\section{Conclusion}

Les résultats obtenus à l'issue du fine-tuning confirment la pertinence du
choix de DistilBERT pour notre cas d'usage. Sur un jeu de validation de
40\,000 exemples totalement indépendant des données d'entraînement, le
modèle retenu (checkpoint-4000) atteint une accuracy de 82,7\% et un F1
de 82,3\%, des performances cohérentes avec l'état de l'art pour un modèle
de cette taille entraîné sur un corpus de taille modeste.

Au-delà des métriques globales, l'analyse des exemples concrets de
prédiction révèle que le modèle distingue correctement et avec une grande
confiance les textes humains courants des textes générés par des LLMs dans
des conditions standards. Ses limites apparaissent principalement sur des
textes générés dans un style littéraire ou rhétorique particulièrement
travaillé — registres moins représentés dans les données d'entraînement —
ce qui constitue une piste d'amélioration identifiée pour des travaux
ultérieurs.

Du point de vue des performances d'inférence, DistilBERT répond pleinement
aux contraintes imposées par le contexte d'une extension navigateur. Sur
une carte graphique grand public RTX~4060 (8\,Go de VRAM), le modèle
produit une prédiction en \textbf{250\,ms} par passage — une latence
imperceptible pour l'utilisateur final, qui rend l'analyse en temps réel
du contenu d'une page web tout à fait réaliste. Ce résultat valide
rétrospectivement le choix de DistilBERT face à des alternatives plus
lourdes comme BERT-base ou RoBERTa, dont le gain de performance marginal
n'aurait pas justifié un temps d'inférence significativement plus élevé
dans ce contexte d'utilisation.

% ════════════════════════════════════════════════════════════════════════════
%              CHAPITRE 5 — SERVICES DU SYSTÈME
% ════════════════════════════════════════════════════════════════════════════

\chapter{Services du système}

\section{Vue d'ensemble}

La partie applicative du projet est organisée sous forme de services séparés.
Chaque service possède une responsabilité limitée : exposition web, routage,
authentification anonyme, limitation de débit, orchestration de la prédiction,
inférence du modèle, stockage ou observabilité.

Cette organisation répond aux principes Cloud-native demandés dans le cahier des
charges. Le système est modulaire, conteneurisé, déployable par environnement et
préparé pour une montée en charge horizontale. Deux environnements sont
structurés dans le projet : un environnement de développement basé sur Docker
Compose et un environnement de staging basé sur Docker Swarm.

Le flux global suit une logique simple : le client récupère d'abord un token
anonyme, l'utilise pour soumettre un ou plusieurs segments textuels, puis reçoit
une réponse normalisée contenant l'identifiant du segment, la décision
\texttt{isAi} et le score associé. Les services internes restent isolés du réseau
public et ne sont accessibles qu'à travers les réseaux Docker dédiés.

\begin{figure}[H]
\centering
\resizebox{\textwidth}{!}{%
\begin{tikzpicture}[
  node distance=5cm and 0.75cm,
  service/.style={draw, rounded corners, align=center, font=\large, inner sep=7pt, minimum width=4.1cm, minimum height=1cm, line width=0.7pt},
  note/.style={align=center, font=\normalsize},
  arrow/.style={-{Stealth[length=3mm]}, very thick}
]
\node[service] (user) {Utilisateur / Extension};
\node[service, below=of user] (proxy) {\texttt{https-proxy}\\Nginx, TLS, routage web et API};
\node[service, below left=of proxy] (frontend) {\texttt{ai-text-detector-frontend}\\Frontend React};
\node[service, below right=of proxy] (gateway) {\texttt{rate-limiting-service}\\Gateway, JWT anonyme, rate limiting};
\node[service, below left=of gateway] (token) {\texttt{token-manager-service}};
\node[service, below right=of gateway] (handler) {\texttt{model-handler-service}};
\node[service, below=of gateway] (redis) {\texttt{redis-server}\\Mémoire partagée};
\node[service, below left=of handler] (postgres) {\texttt{postgres-service}};
\node[service, below right=of handler] (model) {\texttt{transformer-model}\\\texttt{-service}};

\draw[arrow] (user) -- node[right, note]{HTTPS} (proxy);
\draw[arrow] (proxy) -- (frontend);
\draw[arrow] (proxy) -- (gateway);
\draw[arrow] (gateway) -- (token);
\draw[arrow] (gateway) -- (handler);
\draw[arrow] (gateway) -- (redis);
\draw[arrow] (token) -- (redis);
\draw[arrow] (handler) -- (postgres);
\draw[arrow] (handler) -- (model);
\end{tikzpicture}%
}
\caption{Schéma logique simplifié du flux applicatif.}
\label{fig:architecture-logique}
\end{figure}

\section{Synthèse des services}

\begin{table}[H]
\centering
\renewcommand{\arraystretch}{1.35}
\begin{tabular}{p{0.28\textwidth}p{0.55\textwidth}p{0.10\textwidth}}
\toprule
\textbf{Service} & \textbf{Responsabilité} & \textbf{État} \\
\midrule
\texttt{https-proxy} & Reverse proxy Nginx, terminaison HTTPS, routage frontend/API & réalisé \\
\texttt{ai-text-detector}\newline\texttt{-frontend} & Interface React d'analyse manuelle de texte & réalisé \\
\texttt{rate-limiting}\newline\texttt{-service} & API Gateway, contrôle d'accès, vérification JWT, rate limiting & réalisé \\
\texttt{token-manager}\newline\texttt{-service} & Génération de tokens anonymes JWT signés RSA & réalisé \\
\texttt{redis-server} & Mémoire partagée pour les compteurs temporaires & réalisé \\
\texttt{model-handler}\newline\texttt{-service} & Service métier de prédiction, normalisation et historisation & réalisé \\
\texttt{transformer-model}\newline\texttt{-service} & API FastAPI d'inférence du modèle DistilBERT & réalisé \\
\texttt{postgres}\newline\texttt{-service} & Persistance des prédictions et historiques utilisateurs & réalisé \\
\texttt{loki} & Stockage centralisé des logs & réalisé \\
\texttt{promtail} & Collecte et labellisation des logs applicatifs & réalisé \\
\texttt{grafana} & Visualisation, dashboards et suivi opérationnel & réalisé \\
\bottomrule
\end{tabular}
\caption{Services présents dans l'architecture finale.}
\label{tab:services-etat}
\end{table}

\section{Service \texttt{https-proxy}}

Le service \texttt{https-proxy} est le point d'entrée réseau du système. Il est
basé sur Nginx et reçoit le trafic utilisateur. Il termine la connexion HTTPS,
puis route les requêtes vers le frontend ou vers l'API selon leur chemin.

Les routes commençant par \texttt{/api/} sont envoyées vers le
\texttt{rate-limiting-service}. Les autres routes sont envoyées vers le frontend
React. Le proxy ajoute également les headers \texttt{X-Real-IP},
\texttt{X-Forwarded-For} et \texttt{X-Forwarded-Proto}, utilisés par le gateway et
le service d'authentification pour conserver le contexte réseau du client.

\section{Service \texttt{ai-text-detector-frontend}}

Le frontend est une application React/Vite servie en production dans un
conteneur dédié. Il permet à l'utilisateur de saisir un texte, de lancer une
analyse et d'afficher le résultat sous forme de label \texttt{AI} ou
\texttt{HUMAN}, accompagné d'un score de confiance et d'une latence.

Avant d'appeler l'API de détection, le frontend récupère un token anonyme via :

\begin{verbatim}
POST /api/anonym-token/get
\end{verbatim}

Le token est ensuite transmis dans le header :

\begin{verbatim}
Authorization: Anonym <token>
\end{verbatim}

Le frontend conserve aussi un historique local des analyses de session afin
d'afficher des statistiques immédiates : nombre d'analyses, latence moyenne,
longueur des textes et répartition des résultats.

\section{Service \texttt{rate-limiting-service}}

Le \texttt{rate-limiting-service} est le service central de contrôle d'accès.
Il est développé avec Spring Cloud Gateway WebFlux. Il reçoit les appels API
depuis le proxy, applique les politiques d'accès, vérifie les tokens anonymes et
limite le nombre d'appels autorisés.

Deux routes principales sont configurées :

\begin{itemize}
  \item \texttt{/api/anonym-token/get}, route publique contrôlée par la politique
        \texttt{public-token} ;
  \item \texttt{/api/ai-detector/**}, route protégée contrôlée par la politique
        \texttt{protected-anonymous}.
\end{itemize}

Pour les routes protégées, le gateway exige le schéma :

\begin{verbatim}
Authorization: Anonym <token>
\end{verbatim}

Le token est vérifié à l'aide de la clé publique exposée par le
\texttt{token-manager-service} au format JWKS. Le gateway ne possède donc jamais
la clé privée de signature.

Le rate limiting repose sur Redis. Le gateway calcule une empreinte du token,
puis utilise cette empreinte comme clé de compteur. La configuration actuelle
autorise cinq appels par token sur une fenêtre de soixante secondes.

\begin{verbatim}
api-gateway-service.security.ratelimiting.max-attempts-per-token=5
api-gateway-service.security.ratelimiting.window-size-per-token=60
\end{verbatim}

\section{Service \texttt{token-manager-service}}

Le \texttt{token-manager-service} est un service Spring Boot chargé de générer
des tokens anonymes. Il ne gère pas une authentification nominative : son rôle
est de fournir un jeton temporaire permettant d'identifier une session anonyme
et de limiter son usage.

Le service vérifie que la requête provient du gateway grâce à un header secret
partagé. Chaque token contient un identifiant utilisateur anonyme et l'adresse IP
normalisée. Le token est signé avec une clé RSA privée montée dans le conteneur.
La clé publique est exposée via :

\begin{verbatim}
GET /.well-known/jwks.json
\end{verbatim}

Le service limite également le nombre de tokens générés par adresse IP. La
configuration actuelle autorise cinq tokens par IP sur soixante secondes.

\begin{verbatim}
api-gateway-service.security.authentication.token-rate-limit.max-token-per-ip=5
api-gateway-service.security.authentication.token-rate-limit.window-size-per-ip=60
\end{verbatim}

\section{Service \texttt{redis-server}}

Redis sert de mémoire partagée pour les compteurs temporaires. Il est utilisé
par le gateway pour compter les appels par token et par le service
d'authentification pour compter les générations de tokens par IP.

Les données stockées dans Redis sont temporaires :

\begin{itemize}
  \item compteurs de tokens générés par IP ;
  \item compteurs d'appels API par token ;
  \item expiration automatique des clés selon les fenêtres configurées.
\end{itemize}

Redis n'est pas utilisé pour stocker des textes utilisateur ni des résultats de
prédiction.

\section{Service \texttt{model-handler-service}}

Le \texttt{model-handler-service} est le service métier placé entre le gateway
et le service d'inférence. Il reçoit les requêtes applicatives sur :

\begin{verbatim}
POST /api/ai-detector/detect
\end{verbatim}

Le frontend et l'extension envoient une liste d'items textuels. Chaque item
contient au minimum un identifiant et un texte.

\begin{verbatim}
{
  "sourceType": "text",
  "link": null,
  "items": [
    { "id": 1, "text": "Texte à analyser" }
  ]
}
\end{verbatim}

Le modèle retourne des labels de type \texttt{LABEL\_0} et
\texttt{LABEL\_1}. Le service métier transforme ce résultat en une réponse
exploitable par les clients, avec un booléen \texttt{isAi} et un score.

\begin{verbatim}
[
  { "id": 1, "isAi": true, "score": 0.93 }
]
\end{verbatim}

Le service intègre aussi une couche JPA et une entité \texttt{Prediction} afin
de conserver l'historique des analyses associées aux utilisateurs connectés.

\section{Service \texttt{transformer-model-service}}

Le \texttt{transformer-model-service} est le service Python FastAPI qui expose
l'inférence du modèle DistilBERT. Il charge le modèle depuis le dossier monté
\texttt{/workspace/model} et expose l'endpoint interne :

\begin{verbatim}
POST /predict
\end{verbatim}

L'endpoint reçoit une liste de textes :

\begin{verbatim}
[
  { "id": 1, "text": "Texte à analyser" }
]
\end{verbatim}

Il retourne une liste de prédictions :

\begin{verbatim}
[
  { "id": 1, "label": "LABEL_0", "prediction": 0.93 }
]
\end{verbatim}

Ce service est l'un des principaux candidats à la réplication, car l'inférence du
modèle représente l'étape la plus coûteuse en calcul.

\section{Service \texttt{postgres-service}}

PostgreSQL est utilisé par le \texttt{model-handler-service} pour la persistance
des prédictions et des historiques. Le conteneur utilise l'image
\texttt{postgres:16} et un volume Docker nommé \texttt{pgdata}. Les données
conservées permettent de retrouver les analyses associées aux utilisateurs et de
suivre l'évolution des scores moyens par lien ou par session.

\section{Services d'observabilité}

L'observabilité repose sur trois services : Loki, Promtail et Grafana.

Loki centralise les logs applicatifs. Promtail lit les fichiers de logs du
gateway et du service de tokens, puis extrait des labels comme le service, le
niveau, l'environnement, le type d'événement, le statut, l'identifiant de route
et l'état du rate limiting.

Grafana est provisionné automatiquement avec Loki comme source de données. Deux
dashboards principaux sont disponibles :

\begin{itemize}
  \item \texttt{Traffic and Enforcement Overview} ;
  \item \texttt{Security SIEM Overview}.
\end{itemize}

\section{Extension navigateur}

L'extension navigateur est le client principal du système. Elle permet
d'analyser un texte saisi manuellement ou le contenu visible d'une page web. Elle
utilise Manifest V3 et les permissions \texttt{activeTab}, \texttt{scripting} et
\texttt{tabs}.

L'extension parcourt le DOM de la page avec un \texttt{TreeWalker}, ignore les
balises non pertinentes et conserve les segments visibles d'au moins cinquante
caractères. Après l'analyse, elle ajoute temporairement des attributs HTML pour
surligner les segments classés comme IA ou humains.

% ════════════════════════════════════════════════════════════════════════════
%              CHAPITRE 6 — FLUX FONCTIONNELS ET API
% ════════════════════════════════════════════════════════════════════════════

\chapter{Flux fonctionnels et contrats API}

\section{Flux de génération d'un token anonyme}

\begin{figure}[H]
\centering
\fbox{
\begin{minipage}{0.9\textwidth}
\small
\begin{enumerate}
  \item Le client appelle \texttt{POST /api/anonym-token/get}.
  \item Le proxy transmet la requête au gateway avec l'adresse IP client.
  \item Le gateway valide l'IP et ajoute le secret partagé interne.
  \item Le \texttt{token-manager-service} vérifie le secret partagé.
  \item Le service vérifie le compteur Redis de génération par IP.
  \item Le service génère un JWT anonyme signé.
  \item Le token est retourné au client.
\end{enumerate}
\end{minipage}
}
\caption{Flux de génération d'un token anonyme.}
\label{fig:flux-token}
\end{figure}

\section{Flux de prédiction}

\begin{figure}[H]
\centering
\fbox{
\begin{minipage}{0.9\textwidth}
\small
Client $\rightarrow$ \texttt{/api/ai-detector/detect} avec
\texttt{Authorization: Anonym <token>}\\[0.15cm]
$\downarrow$\\
Gateway : vérification JWKS, validation du token, contrôle Redis\\[0.15cm]
$\downarrow$\\
\texttt{model-handler-service} : validation métier et appel du modèle\\[0.15cm]
$\downarrow$\\
\texttt{transformer-model-service} : inférence FastAPI\\[0.15cm]
$\downarrow$\\
Réponse normalisée : \texttt{id}, \texttt{isAi}, \texttt{score}
\end{minipage}
}
\caption{Flux de détection d'un texte.}
\label{fig:flux-detection}
\end{figure}

\section{Endpoints principaux}

\begin{table}[H]
\centering
\small
\renewcommand{\arraystretch}{1.45}
\begin{tabular}{p{0.31\textwidth}p{0.25\textwidth}p{0.34\textwidth}}
\toprule
\textbf{Endpoint} & \textbf{Service cible} & \textbf{Rôle} \\
\midrule
\texttt{POST} \newline \texttt{/api/anonym-token/get} & \texttt{token-manager} \newline \texttt{-service} & Génération d'un token anonyme \\
\texttt{GET} \newline \texttt{/.well-known/jwks.json} & \texttt{token-manager} \newline \texttt{-service} & Publication de la clé publique \\
\texttt{POST} \newline \texttt{/api/ai-detector/detect} & \texttt{model-handler} \newline \texttt{-service} & Analyse d'un ou plusieurs textes \\
\texttt{GET} \newline \texttt{/api/ai-detector/history} & \texttt{model-handler} \newline \texttt{-service} & Historique des analyses connectées \\
\texttt{POST} \newline \texttt{/predict} & \texttt{transformer-model} \newline \texttt{-service} & Inférence interne du modèle \\
\bottomrule
\end{tabular}
\caption{Contrats API principaux.}
\label{tab:endpoints}
\end{table}

\section{Gestion des erreurs}

Les erreurs applicatives sont prises en charge par la couche responsable. Le
gateway retourne une erreur lorsqu'un token est absent, invalide, expiré ou
lorsque la limite de débit est atteinte. Le \texttt{token-manager-service}
retourne une erreur si le secret partagé est incorrect, si l'adresse IP est
absente ou si la limite de génération de tokens est atteinte. Le service métier
retourne une erreur lorsque le payload ne respecte pas le format attendu.

% ════════════════════════════════════════════════════════════════════════════
%              CHAPITRE 7 — CONTENEURISATION ET ORCHESTRATION
% ════════════════════════════════════════════════════════════════════════════

\chapter{Conteneurisation et orchestration par environnement}

\section{Organisation générale}

La conteneurisation du projet repose sur deux niveaux complémentaires. Le dossier
\texttt{runtime-services} contient les définitions unitaires utilisées en
développement. Le dossier \texttt{deploy/staging/services} contient la version
adaptée à l'environnement staging, avec les images du canal \texttt{staging} et
le tag \texttt{v1.0.0}.

Les images applicatives sont publiées sur GitHub Container Registry selon la
convention suivante :

\begin{verbatim}
ghcr.io/baddou666/projet-metier/<channel>/<image>:<tag>
\end{verbatim}

Pour l'environnement staging, les images utilisent le canal
\texttt{staging} et le tag \texttt{v1.0.0}.

\section{Images Docker}

\begin{table}[H]
\centering
\small
\renewcommand{\arraystretch}{1.45}
\begin{tabular}{p{0.30\textwidth}p{0.38\textwidth}p{0.22\textwidth}}
\toprule
\textbf{Service} & \textbf{Image staging} & \textbf{Méthode} \\
\midrule
\texttt{token-manager}\newline\texttt{-service} & \texttt{ghcr.io/.../staging/token-manager:v1.0.0} & Buildpacks Maven \\
\texttt{rate-limiting}\newline\texttt{-service} & \texttt{ghcr.io/.../staging/rate-limiter:v1.0.0} & Buildpacks Maven \\
\texttt{model-handler}\newline\texttt{-service} & \texttt{ghcr.io/.../staging/model-handler:v1.0.0} & Buildpacks Maven \\
\texttt{transformer-model}\newline\texttt{-service} & \texttt{ghcr.io/.../staging/transformer-model:v1.0.0} & Dockerfile \\
\texttt{ai-text-detector}\newline\texttt{-frontend} & \texttt{ghcr.io/.../staging/ai-text-detector-frontend:v1.0.0} & Dockerfile multi-stage \\
\bottomrule
\end{tabular}
\caption{Images applicatives utilisées en staging.}
\label{tab:images-staging}
\end{table}

\section{Environnement de développement}

L'environnement de développement est assemblé par
\texttt{deploy/dev/docker-compose.yml}. Ce fichier inclut les fichiers Compose
unitaires présents dans \texttt{runtime-services}. Cette organisation garde une
configuration par service tout en permettant un lancement global sur une seule
machine.

\section{Environnement staging avec Docker Swarm}

L'environnement staging est orchestré avec Docker Swarm. Il utilise une stack
Swarm dédiée, des réseaux overlay et des images publiées sur le canal
\texttt{staging}. Le lancement se fait avec la commande :

\begin{verbatim}
docker stack deploy -c deploy/staging/docker-compose.yml ai-detect-staging
\end{verbatim}

Le fichier de stack définit les services, les réseaux overlay et les paramètres
de réplication. Les services les plus sollicités, notamment le gateway et le
service d'inférence, sont répliqués afin d'améliorer la disponibilité et la
capacité de traitement.

\section{Réseaux}

Les réseaux staging reprennent la séparation logique du développement, mais sous
forme de réseaux overlay Swarm :

\begin{itemize}
  \item \texttt{proxy-ratelimiter-net} entre Nginx et le gateway ;
  \item \texttt{proxy-fe-net} entre Nginx et le frontend ;
  \item \texttt{redis-net} entre Redis, gateway et token manager ;
  \item \texttt{api-model-net} entre gateway et service métier ;
  \item \texttt{models-net} entre service métier et service d'inférence ;
  \item \texttt{postgres-net} entre service métier et PostgreSQL ;
  \item \texttt{monitoring} pour Loki, Promtail et Grafana.
\end{itemize}

\section{Secrets, volumes et configurations}

Les secrets de staging sont isolés de ceux de développement. Les mots de passe,
clés privées, tokens GHCR et secrets partagés sont fournis via Ansible Vault et
injectés dans l'environnement de déploiement. Les données persistantes sont
stockées dans des volumes Docker nommés, tandis que les fichiers non sensibles
sont fournis comme configurations montées.

% ════════════════════════════════════════════════════════════════════════════
%              CHAPITRE 8 — INFRASTRUCTURE AS CODE
% ════════════════════════════════════════════════════════════════════════════

\chapter{Déploiement et Infrastructure as Code}

\section{Objectif}

L'Infrastructure as Code permet de rendre le déploiement reproductible. Le projet
utilise Terraform pour créer les machines virtuelles et Ansible pour configurer
les noeuds, installer Docker, initialiser Docker Swarm et déployer la stack.

\section{Provisionnement Terraform}

La configuration Terraform de staging crée quatre machines virtuelles destinées
au cluster Docker Swarm :

\begin{itemize}
  \item \textbf{VM1} : manager Swarm ;
  \item \textbf{VM2} : worker backend ;
  \item \textbf{VM3} : worker IA ;
  \item \textbf{VM4} : worker data.
\end{itemize}

Chaque VM est créée à partir d'une image Ubuntu cloud-init, avec accès SSH par
clé publique, hostname dédié et enregistrement Tailscale. Les ressources CPU,
mémoire et disque sont adaptées au rôle du noeud, le noeud IA recevant davantage
de ressources pour supporter l'inférence du modèle.

\section{Configuration Ansible}

Ansible installe les dépendances système, Docker Engine, Docker Compose Plugin,
Git, Git LFS et OpenSSL sur les quatre machines. Le playbook initialise ensuite
Docker Swarm sur le manager, récupère le token de jonction worker, puis joint les
trois autres machines au cluster.

Les noeuds sont labellisés selon leur rôle :

\begin{itemize}
  \item \texttt{role=manager} ;
  \item \texttt{role=backend} ;
  \item \texttt{role=ai} ;
  \item \texttt{role=data}.
\end{itemize}

Ces labels permettent de contrôler le placement des services dans la stack
Swarm.

\section{Déploiement de la stack staging}

Une fois le cluster initialisé, Ansible authentifie les noeuds auprès de GHCR,
récupère le dépôt, prépare les fichiers de configuration et lance la stack
staging avec :

\begin{verbatim}
docker stack deploy -c docker-compose.yml ai-detector-staging
\end{verbatim}

Cette automatisation permet de reconstruire l'environnement staging de manière
reproductible.

\section{Chaîne CI/CD}

La chaîne CI/CD construit les images applicatives, exécute les tests, publie les
images sur GHCR et déclenche le déploiement staging. Les images validées sont
taguées en \texttt{v1.0.0} pour la livraison finale. Cette séparation entre build,
publication et déploiement garantit que le staging utilise des artefacts
identifiés et reproductibles.

% ════════════════════════════════════════════════════════════════════════════
%              CHAPITRE 9 — OBSERVABILITÉ ET SÉCURITÉ
% ════════════════════════════════════════════════════════════════════════════

\chapter{Observabilité et sécurité opérationnelle}

\section{Pipeline d'observabilité}

Les services Spring Boot principaux produisent des logs structurés JSON. Ces logs
contiennent des informations de contexte comme le service, l'environnement, le
type d'événement, l'adresse IP, l'identifiant utilisateur anonyme, la route et le
statut.

\begin{figure}[H]
\centering
\fbox{
\begin{minipage}{0.9\textwidth}
\small
\texttt{rate-limiting-service} et \texttt{token-manager-service}\\
$\downarrow$ logs JSON\\
\texttt{promtail}\\
$\downarrow$ labels et push HTTP\\
\texttt{loki}\\
$\downarrow$ requêtes LogQL\\
\texttt{grafana dashboards}
\end{minipage}
}
\caption{Pipeline d'observabilité basé sur les logs.}
\label{fig:observabilite}
\end{figure}

\section{Dashboards Grafana}

Deux dashboards sont provisionnés :

\begin{itemize}
  \item \texttt{Traffic and Enforcement Overview}, orienté trafic, tokens,
        contrôles et répartition des routes ;
  \item \texttt{Security SIEM Overview}, orienté erreurs, abus, IP bloquées,
        échecs JWKS et événements sensibles.
\end{itemize}

Ces dashboards permettent de suivre le nombre de tokens générés, les requêtes
acceptées, les requêtes refusées, les erreurs de validation de token et les
événements de limitation de débit.

\section{Sécurité applicative}

La sécurité repose sur plusieurs couches :

\begin{itemize}
  \item exposition publique limitée au proxy ;
  \item routage API obligatoire via le gateway ;
  \item secret partagé entre gateway et service de tokens ;
  \item tokens anonymes signés RSA ;
  \item vérification JWKS côté gateway ;
  \item expiration courte des tokens anonymes ;
  \item limitation de génération par IP ;
  \item limitation d'utilisation par token ;
  \item nettoyage des headers transmis aux services internes ;
  \item séparation des secrets par environnement.
\end{itemize}

Ces mécanismes répondent à deux objectifs complémentaires. Le premier est de
protéger les services internes contre les appels directs ou abusifs. Le second
est de limiter la collecte d'informations personnelles : l'utilisateur peut
obtenir un accès temporaire sans créer de compte, tandis que le système conserve
uniquement les informations nécessaires au contrôle de débit et, pour les
utilisateurs connectés, à l'historisation des prédictions.

% ════════════════════════════════════════════════════════════════════════════
%              CHAPITRE 10 — SCALABILITÉ
% ════════════════════════════════════════════════════════════════════════════

\chapter{Scalabilité, réplication et staging}

\section{Principe}

La scalabilité visée est horizontale : augmenter le nombre d'instances de
certains services sans modifier leur code métier. Cette stratégie est rendue
possible par la séparation des responsabilités et l'externalisation de l'état
dans Redis et PostgreSQL.

\section{Réplication du \texttt{rate-limiting-service}}

Le gateway est réplicable car son état critique de rate limiting est externalisé
dans Redis. Plusieurs instances peuvent partager les mêmes compteurs, à
condition que toutes utilisent la même instance Redis et la même configuration de
routes, de secrets et de JWKS.

\section{Réplication du \texttt{transformer-model-service}}

Le service d'inférence est réplicable car il ne stocke pas de session utilisateur
locale. Chaque instance charge le modèle depuis son dossier
\texttt{/workspace/model} et peut traiter une requête indépendamment. La
réplication permet de répartir la charge liée à l'inférence.

\section{Placement des services}

Dans le cluster Docker Swarm, les noeuds sont séparés par rôle :

\begin{itemize}
  \item le manager reçoit les services d'orchestration et le proxy ;
  \item le worker backend reçoit le gateway, le token manager et le model handler ;
  \item le worker IA reçoit les instances du service d'inférence ;
  \item le worker data reçoit Redis, PostgreSQL, Loki et Grafana.
\end{itemize}

Cette répartition permet d'isoler les charges de calcul, les charges réseau et
les charges de stockage.

\section{Robustesse de la réplication}

La réplication horizontale améliore la disponibilité et le débit du système en
répartissant les appels sur plusieurs instances applicatives. Les services
répliqués restent stateless : le gateway délègue ses compteurs à Redis, tandis
que le service d'inférence traite chaque requête indépendamment. Les composants
à état, notamment Redis et PostgreSQL, sont isolés sur le worker data avec des
volumes persistants et une supervision dédiée.

Cette organisation permet de stabiliser les charges : les requêtes entrantes
sont absorbées par les instances du gateway, les opérations de prédiction sont
réparties entre les instances du modèle, et les données persistantes restent
centralisées dans les services spécialisés. Le staging Docker Swarm valide ainsi
la séparation entre capacité de calcul, stockage et routage réseau.

% ════════════════════════════════════════════════════════════════════════════
%              CHAPITRE 11 — VALIDATION
% ════════════════════════════════════════════════════════════════════════════

\chapter{Validation et tests}

\section{Validation fonctionnelle}

Le flux nominal validé est le suivant :

\begin{enumerate}
  \item démarrer la stack Docker Compose ou Docker Swarm ;
  \item appeler \texttt{POST /api/anonym-token/get} via le proxy ;
  \item récupérer le JWT anonyme ;
  \item appeler \texttt{POST /api/ai-detector/detect} avec le header
        \texttt{Authorization: Anonym <token>} ;
  \item vérifier la réponse \texttt{id}, \texttt{isAi}, \texttt{score} ;
  \item vérifier les logs et les dashboards dans Grafana.
\end{enumerate}

\section{Tests d'intégration}

Les tests d'intégration couvrent le passage complet entre le proxy, le gateway,
le token manager, Redis, le model handler et le service FastAPI. Ils permettent
de vérifier la génération de token, la validation JWKS, le rate limiting et le
format final des réponses de prédiction.

\begin{table}[H]
\centering
\renewcommand{\arraystretch}{1.35}
\begin{tabular}{p{0.36\textwidth}p{0.45\textwidth}p{0.10\textwidth}}
\toprule
\textbf{Scénario} & \textbf{Vérification attendue} & \textbf{État} \\
\midrule
Génération de token anonyme & Retour d'un JWT signé et exploitable par le gateway & validé \\
Appel sans token & Rejet de la requête protégée & validé \\
Appel avec token invalide & Rejet après vérification JWKS & validé \\
Limitation de débit & Blocage après dépassement du seuil configuré & validé \\
Prédiction nominale & Retour de \texttt{id}, \texttt{isAi} et \texttt{score} & validé \\
Collecte des logs & Présence des événements dans Loki et Grafana & validé \\
\bottomrule
\end{tabular}
\caption{Synthèse des tests d'intégration réalisés.}
\label{tab:tests-integration}
\end{table}

\section{Tests de charge}

Des tests de charge ont été réalisés sur les routes principales afin de valider
la réplication du gateway et du service d'inférence. Les résultats confirment que
le service \texttt{transformer-model-service} est le principal facteur de coût,
ce qui justifie son placement sur le worker IA et sa réplication dans la stack
staging.

\section{Bilan de validation}

La validation confirme que le système répond au flux utilisateur visé : analyse
d'un contenu textuel depuis un client web ou une extension, passage sécurisé par
le gateway, prédiction par le modèle et restitution d'un score exploitable. Les
tests montrent également que les mécanismes transverses, notamment la limitation
de débit, la vérification de token et l'observabilité, sont intégrés au flux
principal et ne restent pas des composants isolés.

% ════════════════════════════════════════════════════════════════════════════
%              CHAPITRE 12 — CONCLUSION
% ════════════════════════════════════════════════════════════════════════════

\chapter{Conclusion générale}

Le système réalisé met en place une architecture microservices complète autour
d'un détecteur de textes générés par IA. Le coeur livré comprend un frontend
React, une extension navigateur, un proxy HTTPS, un gateway avec contrôle
d'accès, un service de tokens anonymes, un service métier de prédiction, un
service FastAPI d'inférence, Redis, PostgreSQL et une pile d'observabilité
Loki/Promtail/Grafana.

L'architecture respecte les exigences principales du cahier des charges :
conteneurisation, séparation des responsabilités, API de prédiction,
orchestration par environnement, observabilité, limitation de débit,
reproductibilité et scalabilité horizontale.

Le déploiement final est organisé autour de deux environnements : un
environnement de développement basé sur Docker Compose et un environnement
staging basé sur Docker Swarm. L'Infrastructure as Code, construite avec
Terraform et Ansible, permet de provisionner les machines, de configurer Docker,
d'initialiser le cluster Swarm et de déployer la stack de manière reproductible.

La solution obtenue ne se limite donc pas à un modèle de classification : elle
constitue un système complet, exploitable et conforme à une démarche
Cloud-native, depuis la préparation des données jusqu'au déploiement et au suivi
opérationnel.

Sur le plan expérimental, le choix de DistilBERT s'est révélé adapté aux
contraintes du projet. Le modèle conserve une latence faible tout en atteignant
des performances suffisantes pour assister l'utilisateur dans l'identification
de contenus potentiellement générés par IA. Le résultat ne doit toutefois pas
être interprété comme une preuve absolue d'origine artificielle : il s'agit d'un
indicateur probabiliste, sensible au style, au domaine et aux transformations
éventuelles du texte.

Le projet répond ainsi au cahier des charges initial : ingestion et préparation
des données, entraînement d'un modèle NLP avancé, exposition d'une API de
prédiction, déploiement conteneurisé, orchestration multi-environnement,
scalabilité, sécurité d'accès et observabilité. La version finale fournit une
chaîne complète allant du texte brut jusqu'à l'analyse visuelle dans le
navigateur, avec une infrastructure reproductible et un suivi opérationnel
adapté à une exploitation Cloud-native.

\end{document}

